% SEGMENT OLP-0021-S01
% Part: sets-functions-relations
% Chapter: functions
% Section: basics

\documentclass[../../../include/open-logic-section]{subfiles}

\begin{document}

\olfileid{sfr}{fun}{bas}
\olsection{அடிப்படைகள்}

\begin{explain}
ஒரு \emph{சார்பு} என்பது, கொடுக்கப்பட்ட ஒரு கணத்தின் ஒவ்வொரு
!!{element}[acc]யும், கொடுக்கப்பட்ட ஏதோ ஒரு (வேறு) கணத்தின் குறிப்பிட்ட
!!{element} ஒன்றுடன் இணைக்கும் ஒரு இணைப்பு. எடுத்துக்காட்டாக,
$1$ ஐக் கூட்டும் செயல் ஒரு சார்பை வரையறுக்கிறது: ஒவ்வொரு எண் $n$ உம்
ஒரே ஒரு எண் $n+1$ உடன் இணைக்கப்படுகிறது.

இன்னும் பொதுவாக, சார்புகள் ஜோடிகள், மூன்று-தொகுதிகள் முதலியவற்றை
உள்ளீடுகளாகப் பெற்று ஏதோ ஒரு வகையான வெளியீட்டைத் தரலாம்.
அடிப்படை எண்கணிதத்திலிருந்து பல சார்புகளை நாம் அறிவோம்.
எடுத்துக்காட்டாக, கூட்டலும் பெருக்கலும் சார்புகளே. அவை இரண்டு எண்களைப்
பெற்று மூன்றாவது ஓர் எண்ணைத் தருகின்றன.

இந்த அருவமான கணிதப் பொருளில், ஒரு சார்பு ஒரு \emph{கருப்புப் பெட்டி}:
எந்த உள்ளீட்டுடன் எந்த வெளியீடு இணைக்கப்படுகிறது என்பதே முக்கியம்;
வெளியீட்டைக் கணக்கிடும் முறை முக்கியமல்ல.
\end{explain}

% SEGMENT OLP-0021-S02
\begin{defn}[சார்பு]
ஒரு \emph{சார்பு} $f \colon A \to B$ என்பது $A$ இன் ஒவ்வொரு
!!{element}[acc]யும் $B$ இன் ஓர் !!{element} உடன் இணைக்கும் இணைப்பாகும்.

$A$ ஐ $f$ இன் \emph{சார்பகம்} என்றும், $B$ ஐ $f$ இன்
\emph{துணைச்சார்பகம்} என்றும் அழைக்கிறோம்.
$A$ இன் !!{element}s, $f$ இன் உள்ளீடுகள் அல்லது
\emph{உள்ளீட்டுறுப்புகள்} எனப்படும். ஓர் உள்ளீட்டுறுப்பு $x$ உடன்
$f$ இணைக்கும் $B$ இன் !!{element}, உள்ளீட்டுறுப்பு $x$ க்கான
\emph{$f$ இன் மதிப்பு} எனப்படும்; அது $f(x)$ என எழுதப்படுகிறது.

$f$ இன் \emph{வீச்சகம்} $\ran{f}$ என்பது, ஏதாவது ஓர்
உள்ளீட்டுறுப்புக்கான $f$ இன் மதிப்புகளைக் கொண்ட துணைச்சார்பகத்தின்
உட்கணமாகும்; $\ran{f} =
\Setabs{f(x)}{x \in A}$.
\end{defn}

% SEGMENT OLP-0021-S03
\olref{fig:function} இல் உள்ள படம் சார்புகளைப் பற்றிச் சிந்திக்க உதவலாம்.
இடப்புற நீள்வட்டம் சார்பின் \emph{சார்பகத்தைக்} குறிக்கிறது;
வலப்புற நீள்வட்டம் அதன் \emph{துணைச்சார்பகத்தைக்} குறிக்கிறது.
சார்பகத்தில் உள்ள ஓர் \emph{உள்ளீட்டுறுப்பிலிருந்து}, துணைச்சார்பகத்தில்
அதற்குரிய \emph{மதிப்பை} நோக்கி ஓர் அம்பு செல்கிறது.

\begin{figure}
  \olasset{assets/diagrams/function.tikz}
  \caption{ஒரு சார்பு என்பது ஒரு கணத்தின் ஒவ்வொரு !!{element}[acc]யும்
    மற்றொரு கணத்தின் !!a{element} உடன் இணைக்கும் இணைப்பாகும்.
    சார்பகத்தில் உள்ள ஓர் உள்ளீட்டுறுப்பிலிருந்து, துணைச்சார்பகத்தில்
    அதற்குரிய மதிப்பை நோக்கி ஓர் அம்பு செல்கிறது.}
  \ollabel{fig:function}
\end{figure}

% SEGMENT OLP-0021-S04
\begin{ex}
பெருக்கல் இயல் எண்களின் ஜோடிகளை உள்ளீடுகளாகப் பெற்று, இயல் எண்களை
வெளியீடுகளாகத் தருகிறது. எனவே அது $\Nat \times \Nat$ என்ற சார்பகத்திலிருந்து
$\Nat$ என்ற துணைச்சார்பகத்திற்குச் செல்கிறது.
வீச்சகமும் $\Nat$ ஆகவே அமைகிறது. ஏனெனில் ஒவ்வொரு $n \in \Nat$ உம்
$n \times 1$ ஆகும்.
\end{ex}

% SEGMENT OLP-0021-S05
\begin{ex}
பெருக்கல் ஒவ்வொரு உள்ளீட்டையும், அதாவது இயல் எண்களின் ஒவ்வொரு ஜோடியையும்,
ஒரே ஒரு வெளியீட்டுடன் இணைப்பதால் அது ஒரு சார்பு:
$\times \colon \Nat^2 \to
\Nat$. ஆனால் ஓர் இயல் எண் $n$ ஐ, $x^2 = n$ என அமையும் மெய்யெண்
$x$ உடன் இணைப்பது சார்பாக அமையாது. ஏனெனில் ஒவ்வொரு நேர்ம முழு $n$ க்கும்
இரண்டு வர்க்க மூலங்கள் உள்ளன: $\sqrt{n}$, $-\sqrt{n}$.
எதிர்மமற்ற முதன்மை வர்க்க மூலத்தை மட்டும் திருப்பித் தருவதன் மூலம் இதைச் சார்பாக
மாற்றலாம்: $\sqrt{\phantom{X}} \colon
\Nat \to \Real$.
\end{ex}

% SEGMENT OLP-0021-S06
\begin{ex}
ஒரு வகுப்பிலுள்ள ஒவ்வொரு மாணவரையும் அவருடைய இறுதி மதிப்பெண்ணுடன்
இணைக்கும் தொடர்பு ஒரு சார்பாகும். ஒரே வகுப்பில் எந்த மாணவருக்கும்
இரண்டு வெவ்வேறு இறுதி மதிப்பெண்கள் கிடைக்க முடியாது.
ஒரு வகுப்பிலுள்ள ஒவ்வொரு மாணவரையும் அவருடைய பெற்றோருடன் இணைக்கும்
தொடர்பு சார்பல்ல: மாணவர்களுக்குப் பெற்றோர் இல்லாமலும் இருக்கலாம்;
இருவர் அல்லது அதற்கு மேற்பட்டவர்கள் இருக்கலாம்.
\end{ex}

% SEGMENT OLP-0021-S07
\begin{explain}
சாத்தியமான ஒவ்வொரு உள்ளீட்டுறுப்புக்கும் சார்பின் மதிப்பு என்ன என்பதை
ஏதாவது ஒரு துல்லியமான முறையில் குறிப்பிடுவதன் மூலம் சார்புகளை
வரையறுக்கலாம். ஒரு வாய்பாட்டைத் தருவது, மதிப்பைக் கணக்கிடும் முறையை
விவரிப்பது, அல்லது ஒவ்வொரு உள்ளீட்டுறுப்புக்கும் உரிய மதிப்பைப்
பட்டியலிடுவது ஆகியவை இதற்கான வெவ்வேறு வழிகள். எந்த முறையில் சார்பை
வரையறுத்தாலும், ஒவ்வொரு உள்ளீட்டுறுப்புக்கும் ஒரே ஒரு மதிப்பை மட்டுமே
குறிப்பிடுகிறோம் என்பதை உறுதிசெய்ய வேண்டும்.
\end{explain}

% SEGMENT OLP-0021-S08
\begin{ex}
$f \colon \Nat \to \Nat$ என்பதை $f(x) = x+1$ என வரையறுப்போம்.
இந்த வரையறை $f$ ஐ, இயல் எண்களை உள்ளீடுகளாகப் பெற்று இயல் எண்களை
வெளியீடுகளாகத் தரும் சார்பாகக் குறிப்பிடுகிறது.
ஓர் இயல் எண் $x$ தரப்பட்டால், $f$ அதன் தொடரி $x+1$ ஐ வெளியீடாகத்
தரும் என இது கூறுகிறது. இங்கு துணைச்சார்பகம் $\Nat$ என்பது $f$ இன்
வீச்சகம் அல்ல. ஏனெனில் இயல் எண் $0$ எந்த இயல் எண்ணின் தொடரியும் அல்ல.
$f$ இன் வீச்சகம் அனைத்து நேர்ம முழுக்களின் கணமான $\Int^{+}$ ஆகும்.
\end{ex}

% SEGMENT OLP-0021-S09
\begin{ex}\ollabel{examplefunext}
$g \colon \Nat \to \Nat$ என்பதை $g(x) = x+2-1$ என வரையறுப்போம்.
$g$ இயல் எண்களை உள்ளீடுகளாகப் பெற்று இயல் எண்களை வெளியீடுகளாகத் தரும்
சார்பு என இது கூறுகிறது. ஓர் இயல் எண் உள்ளீடாகத் தரப்பட்டால், $g$ என்பது
$x$ இன் தொடரியின் தொடரிக்கு உரிய முன்னியை, அதாவது $x+1$ ஐ,
வெளியீடாகத் தரும்.
\end{ex}

% SEGMENT OLP-0021-S10
\begin{explain}
வெவ்வேறு \emph{வரையறைகளைக்} கொண்ட $f$, $g$ என்ற இரு சார்புகளைக்
கருதினோம். எனினும் இவை \emph{ஒரே சார்புதான்}.
ஏனெனில் எந்த இயல் எண் $n$ க்கும் $f(n) = n+1 = n+2-1 =
g(n)$. வேறு சொற்களில், $f$, $g$ ஆகியவற்றுக்கான நமது வரையறைகள்
வெவ்வேறு சமன்பாடுகளின் மூலம் ஒரே இணைப்பைக் குறிப்பிடுகின்றன.
எனவே சார்புகளுக்கான மதிப்புசார் சமத்துவக் கொள்கையை மறைமுகமாகப்
பயன்படுத்துகிறோம்:
\[
  \forall x\, f(x) = g(x)\text{ எனில் }f = g
\]
இங்கு $f$, $g$ ஆகியவற்றுக்கு ஒரே சார்பகமும் ஒரே துணைச்சார்பகமும்
இருக்க வேண்டும்.
\end{explain}

% SEGMENT OLP-0021-S11
\begin{ex}
நிலைகளைப் பிரித்தும் சார்புகளை வரையறுக்கலாம். எடுத்துக்காட்டாக,
$h \colon \Nat \to \Nat$ என்பதைப் பின்வருமாறு வரையறுக்கலாம்:
\[
h(x) =
\begin{cases}
  \frac{x}{2} & \text{$x$ இரட்டை எனில்} \\
  \frac{x+1}{2} & \text{$x$ ஒற்றை எனில்.}
\end{cases}
\]
ஒவ்வொரு இயல் எண்ணும் இரட்டையாகவோ ஒற்றையாகவோ இருப்பதால், இந்தச்
சார்பின் வெளியீடு எப்போதும் ஓர் இயல் எண்ணாக இருக்கும்.
நிலைகளைப் பிரித்து ஒரு சார்பை வரையறுக்கும்போது, சாத்தியமான ஒவ்வொரு
உள்ளீடும் சரியாக ஒரு நிலையில் இடம்பெற வேண்டும் என்பதை நினைவில் கொள்க.
சில சமயங்களில், நிலைகள் அனைத்துச் சாத்தியங்களையும் உள்ளடக்குகின்றன
என்பதையும் ஒன்றுக்கொன்று விலக்கானவை என்பதையும் நிறுவ வேண்டியிருக்கும்.
\end{ex}

\end{document}
